% Reproducibility appendix for the BISON manuscript.
% Author: fangtoast
% Email: fangtoast@foxmail.com
% Copyright (c) 2026 fangtoast. All rights reserved.

\section{Reproducibility Details}

\subsection{BISON Configuration}

\begin{table}[H]
  \caption{Architecture and training configuration used for the formal forward
  BISON result. Output dimensions are supplied by the dataset schema.}
  \label{tab:bison-config}
  \centering
  \small
  \begin{tabular}{ll}
    \toprule
    Item & Value \\
    \midrule
    Input representation & periodic permittivity grid \\
    Lift width / BandBlocks & 32 / 4 \\
    Spectral modes / band modes & $8\times8$ / $5\times5$ \\
    Global low-rank width & 8 \\
    Activation & GELU \\
    Pooling / head & adaptive average / linear \\
    Output in this dataset & $31\times10$ bands \\
    Trainable parameters & 389,718 \\
    Selected learning rate & $10^{-3}$ \\
    Seeds & 0, 1, 2 \\
    Evaluation bands & 2--10 \\
    \bottomrule
  \end{tabular}
\end{table}

\subsection{Per-Band Forward Error}

\begin{table}[H]
  \caption{BISON test MRE by evaluated band, averaged over three seeds.}
  \label{tab:per-band}
  \centering
  \small
  \resizebox{\columnwidth}{!}{%
  \begin{tabular}{lrrrrrrrrr}
    \toprule
    Band & 2 & 3 & 4 & 5 & 6 & 7 & 8 & 9 & 10 \\
    \midrule
    MRE (\%) & 1.020 & 0.916 & 0.903 & 0.994 & 1.002 & 0.990 & 1.024 & 1.144 & 1.276 \\
    \bottomrule
  \end{tabular}}
\end{table}

\subsection{Inverse Candidate and Replay Settings}

\begin{table}[H]
  \caption{Frozen seed-0 inverse protocol.}
  \label{tab:inverse-config}
  \centering
  \small
  \begin{tabular}{ll}
    \toprule
    Item & Value \\
    \midrule
    Conditional generator & one fixed mode, no routing \\
    Latent / hidden / depth & 32 / 256 / 4 \\
    Objective fine-tuning & 50 epochs, AdamW, $10^{-4}$ \\
    Candidates per target & 128 \\
    Objective candidates in training & 16 \\
    Binary threshold & 0.5 \\
    Primary ranker & frozen seed-0 BISON \\
    Cross-surrogate diagnostic & frozen seed-0 FNO query model \\
    MPB replay budget & Top 5 per target \\
    Test targets & complete 625-sample split \\
    Bootstrap samples & 10,000 \\
    \bottomrule
  \end{tabular}
\end{table}

\subsection{Artifact Provenance and Claim Status}
\label{app:provenance}

The source manifest retains complete paths and digests. Its SHA-256 prefixes
are listed separately to preserve readable two-column typesetting:
\begin{itemize}
  \item dataset: \texttt{991cb3810e88};
  \item complete-test aggregate: \texttt{af97717abbd5};
  \item per-target replay table: \texttt{716bb6f426c4};
  \item paired qualitative inverse figure: \texttt{62830e47306e}.
\end{itemize}

The forward result is \emph{formal completed} with three seeds. The inverse
result is \emph{single-seed complete-test replay}: it covers all 625 held-out
targets but uses only seed 0. The experiment does not authorize claims of
cross-seed inverse stability, universal superiority, PINO/field inverse gain,
or causal attribution against an unmatched historical seed-2 run.

\section{Additional Interpretation}

The single-mode generator's MPB MRE interval is [2.2263, 2.2965]\%, whereas
the four-mode interval is [2.3485, 2.4341]\%. The paired interval reported in
the main paper is more informative for the same targets because it removes much
of the target-to-target difficulty variation. Even so, it remains conditional
on one trained pair of inverse models. A stronger follow-up would repeat the
entire objective-bridge training and replay protocol for multiple seeds while
preserving the same target pool and MPB budget.
